\begin{afterword}
\summaryhead

The essay argues that people become good reasoners only when something they care about more than reasoning is at stake. It starts from a contrast in fiction: in Western comics heroes get powers first and then look for a purpose, while in Japanese fiction the power comes from having someone to protect.

It then claims that science defeated authority because it produced useful technology, not because it sounded reasonable. It answers the objection that truth should come before usefulness by saying the relationship is complicated. Its main example is a choice between saving 400 people for certain and saving 500 with probability 90\%. A parent whose daughter is among the 500 is more likely, the author suggests, to do the arithmetic and choose the gamble.

The author dates the start of real growth as a reasoner to having something important to do. The essay concludes that one should find a truly important cause by learning to calculate, and put reasoning in its service.

\responsehead

The essay's claim is strong: no one masters rationality until something more than their own life is at stake. Its support is fiction, which it disclaims as evidence; a sketch of the history of science; one numerical dilemma; and the author's own history. None of these supports the claim.

The fiction contains an error that matters to the contrast. The Western hero bitten by a radioactive spider does not fight crime to keep busy. Spider-Man lets a burglar escape, the burglar kills his uncle, and he fights crime out of guilt. His motive is the loss of someone he loved, which is the Japanese pattern the essay sets against him.

The history is asserted in a few sentences: science won ``because it displayed greater raw strength in the form of technology.'' No historian is cited, and the one study I checked complicates it. Britain's great inventors of the industrial period were mostly not scientists until late in the nineteenth century.

The dilemma is the essay's one argument, and its outcome is fixed by arithmetic. From the daughter's side, option 1 saves her with probability 80\% and option 2 with 90\%. Neither is certain, so the ``comforting feeling of certainty'' that the essay says love will overcome has already been removed by the framing. And 80 against 90 is the same ratio as 400 against 450. While the parent does not know which of the 500 is the daughter, love for one child and counting everyone must give the same answer, since her chance is the expected number saved divided by 500. The case cannot show love correcting a calculation, because here the two never differ. What evidence exists on emotion and probability suggests that strong feeling makes people less sensitive to differences like 80 and 90, not more.

That leaves the author's history, which is one person's self-report. The author grew up an ``above-average Traditional Rationalist \ldots\ and nothing more,'' and grew further only once there was ``something terribly important'' to do. The essay does not say what it was. A later post says the author was raised ``in technophilia to treasure that one most precious thing, far more important than my own life.'' My inference, which the essay does not state, is that the cause is the author's own work on artificial intelligence. The essay's advice to readers is to ``Learn how to multiply, and perhaps you will recognize a drastically important cause when you see one.'' Read with that inference, the cause a reader is meant to recognize is the author's.

In short: a universal claim about how people learn to reason, supported by a misremembered comic, an unsourced history, a dilemma whose two answers agree by arithmetic, and one autobiography whose central cause goes unnamed.
\end{afterword}
